\documentclass[12pt,a4paper]{article}

% ---------- Packages ----------
\usepackage[utf8]{inputenc}
\usepackage[T1]{fontenc}
\usepackage[english]{babel}
\usepackage{lmodern}
\usepackage[margin=2.5cm,headheight=15pt]{geometry}
\usepackage{fancyhdr}
\usepackage{lastpage}
\usepackage{graphicx}
\usepackage{float}
\usepackage{booktabs}
\usepackage{tabularx}
\usepackage{xcolor}
\usepackage{listings}
\usepackage{caption}
\usepackage[hidelinks]{hyperref}
\usepackage[nameinlink]{cleveref}

\graphicspath{{figures/}}

% ---------- Your details (the exam wants the full name on EVERY page) ----------
\newcommand{\fullname}{FULL NAME}
\newcommand{\username}{LTU-USERNAME}

% ---------- Header / footer: name on every page, pages numbered ----------
\pagestyle{fancy}
\fancyhf{}
\lhead{\fullname\ (\username)}
\rhead{D7032E Home Exam 2026}
\cfoot{Page \thepage\ of \pageref{LastPage}}

% ---------- Code listings ----------
\lstset{
  language=Java,
  basicstyle=\ttfamily\small,
  keywordstyle=\color{blue}\bfseries,
  commentstyle=\color{green!40!black},
  numbers=left,
  numberstyle=\tiny\color{gray},
  breaklines=true,
  frame=single,
  showstringspaces=false,
  captionpos=b
}

% ---------- Helpers ----------
\newcommand{\todo}[1]{\textcolor{red}{\textbf{TODO:} #1}}   % delete every \todo before hand-in
\newcommand{\req}[1]{\textbf{R#1}}                          % \req{3} -> R3
\newcommand{\code}[1]{\texttt{#1}}

\begin{document}

\begin{center}
  {\Large\bfseries Home Exam -- Software Engineering (D7032E) -- 2026}\\[0.4em]
  {\large Exploding Kittens}\\[0.4em]
  \fullname\ (\username)
\end{center}

% =====================================================================
\section*{Part A}
% =====================================================================

\section{Unit testing (2p, max 1 page)}
\label{sec:q1}

\todo{Running text. Go through requirements 1--13 against the ORIGINAL code
(\code{ExplodingKittens.java}) and name every requirement that is not fulfilled,
using the requirement number. Start with the deck set-up around the \code{initGame}
method and then the turn loop in \code{game}.}

\todo{For each unfulfilled requirement, either give the JUnit assertion that would
detect it WITHOUT modifying the existing code, or motivate why that is impossible
(think: static state, the \code{System.exit} call, I/O inside the logic, no way to
inject the deck or the random source). One short paragraph per requirement:}

\paragraph{\req{X} -- short name.}
\todo{Why it is not fulfilled.}
\begin{lstlisting}[caption={Possible JUnit check for \req{X}.}]
// TODO: assertEquals(expected, actual);
\end{lstlisting}

% =====================================================================
\section{Software architecture design and refactoring (9p, max 2 pages excluding diagrams)}
\label{sec:q2}

\todo{The oral exam focuses on this section. Prefer clear diagrams over long text;
the reasoning can be given verbally. Keep the text within 2 pages, diagrams excluded.
Cite external sources with \code{\textbackslash cite\{\}}.}

\subsection{Shortcomings of the original design (1p)}
\label{sec:q2a}

The original design is assessed with the SOLID principles~\cite{martin2003} and
Booch's metrics~\cite{booch2007} (coupling, cohesion, sufficiency, completeness and
primitiveness).
\todo{One or two concrete shortcomings per SOLID principle and per Booch metric, each
tied to something visible in the code (class, method or line). Keep it short.}

\subsection{Design choices in the new design (3p)}
\label{sec:q2b}

\todo{Overview of the architecture, then the reasoning behind the choices in terms of
SOLID and Booch's metrics. The documentation must let a developer see where each
functionality lives, what the interfaces are, how classes and modules relate and how
communication flows. Diagrams:}

\begin{figure}[H]
  \centering
  % \includegraphics[width=\textwidth]{package-diagram.pdf}
  \fbox{\parbox[c][5cm][c]{0.9\textwidth}{\centering \todo{Package / component diagram}}}
  \caption{Packages and their dependencies.}
  \label{fig:packages}
\end{figure}

\begin{figure}[H]
  \centering
  % \includegraphics[width=\textwidth]{class-diagram.pdf}
  \fbox{\parbox[c][7cm][c]{0.9\textwidth}{\centering \todo{UML class diagram: key classes, interfaces, relations}}}
  \caption{Class diagram of the refactored design.}
  \label{fig:classes}
\end{figure}

\begin{figure}[H]
  \centering
  % \includegraphics[width=\textwidth]{sequence-diagram.pdf}
  \fbox{\parbox[c][6cm][c]{0.9\textwidth}{\centering \todo{Sequence diagram: communication flow, e.g.\ a card played and Noped}}}
  \caption{Communication flow between client, network and game engine.}
  \label{fig:sequence}
\end{figure}

\subsection{Quality attributes (3p)}
\label{sec:q2c}

\todo{One separate subsection per quality attribute. For each: which specific design
choices satisfy it, and which requirement (14--20) it serves.}

\subsubsection{Extensibility}
\todo{New expansions (R14) without changing existing, compiled code.}

\subsubsection{Modifiability}
\todo{Core rules affected by expansion rules (R14); menu (R15); bots (R16); decoupled
network layer (R17).}

\subsubsection{Testability}
\todo{Dependency injection of deck, random source, players and I/O (R19).}

\subsubsection{Concurrency (bonus)}
\todo{Nope window and interrupts (R20); the original code around lines 148 and 163.}

\subsection{Design patterns (2p)}
\label{sec:q2d}

\todo{For each pattern used: the problem it solves here and how it improves the design.
Also say which obvious patterns you chose NOT to use and why. Cite
\cite{gamma1994} where relevant.}

% =====================================================================
\section*{Part B}
% =====================================================================

\section{Quality attributes, design patterns, documentation, re-engineering, testing (14p)}
\label{sec:q3}

\todo{Evaluate the delivered code. Be honest and concrete: what works, what does not,
and evidence (test results, what you had to change to add something). Running text,
short.}

\subsection{Extensibility (2p)}
\todo{To what degree can expansions be added without affecting the existing (compiled)
implementation, including unit- and integration-testing an expansion before launch?
Example: what would adding one Imploding Kittens card require?}

\subsection{Modifiability (3p)}
\todo{To what degree can the code be modified: core rules affected by expansions, and
the network implementation?}

\subsection{Testability (1p)}
\todo{To what degree is the code designed for testability?}

\subsection{Unit-test coverage of requirements 1--5 (2p)}
\todo{To what degree is the code unit-tested? Relate tests to requirements; give the
command to run them (see \cref{app:run}).}

\subsection{Best practices (1p)}
\todo{Structure, standards, naming.}

\subsection{Correctness of game functionality (2p)}
\todo{To what degree does the game implement requirements 1--13 correctly? Known gaps.}

\subsection{Error and exception handling (1p)}
\todo{Are errors and exceptions handled and reported appropriately?}

\subsection{Documentation (1p)}
\todo{Is the code appropriately documented?}

\subsection{Consistency with the design (1p)}
\todo{Is the code true to the design in \cref{sec:q2}? Note any deviations and update
the design if needed.}

% =====================================================================
\section*{Use of AI tools}
% =====================================================================
\todo{State which AI tools you used and for what. You must understand and be able to
justify everything you hand in.}

% =====================================================================
\appendix
\section{How to compile, run and test}
\label{app:run}
\todo{Keep in sync with README.md. Required by the exam: instructions to compile and
run the code and to run the unit tests.}

\begin{lstlisting}[language={},caption={Build, test and run.}]
./mvnw test                                     # run the unit tests
./mvnw package                                  # build target/exploding-kittens.jar
java -jar target/exploding-kittens.jar 2 0      # server: players, bots
java -jar target/exploding-kittens.jar 127.0.0.1  # client: server IP
\end{lstlisting}

% =====================================================================
% References: every source listed here must be cited in the text above.
\bibliographystyle{plain}
\bibliography{references}

\end{document}
